% Tracce uscite agli Esami di Stato (Sezione A, Settore Informazione) sugli argomenti del capitolo.
% Incluso da main.tex subito dopo il riassunto.

\section*{Tracce Uscite} % Usiamo * per non numerare la sezione nell'indice
\addcontentsline{toc}{section}{Tracce Uscite: Progettazione di Sistemi}

\subsection*{Prima sessione 2026 -- A3-INF-4 (prova pratica di progettazione)}
\textbf{Traccia}: cinema multisala con interfaccia Web, backend, base di dati relazionale e modulo di Machine Learning. Nove punti in tre parti: (1) ER, (2) modello relazionale, (3) SQL e query dei posti disponibili; (4) componenti del backend, (5) API REST, (6) diagramma UML di sequenza per la prenotazione con prevenzione dei doppioni; (7) progetto del modello di previsione, (8) addestramento e valutazione, (9) uso delle predizioni e limiti.

\textbf{Dove trovare la risposta}: sezione \emph{Tema Svolto: Cinema Multisala}, che risponde a ciascun punto con la stessa numerazione.

\subsection*{Prima sessione 2018 -- Informatica (seconda prova)}
\textbf{Traccia}: progettare un algoritmo che confronta due immagini $A$ e $B$ di dimensione $W \times H$ per trovare le corrispondenze tra pixel (ricerca del pixel più simile in $B$ nell'intervallo $[x-d]$, $d \in [0,D]$), restituendo una nuova immagine con lo spostamento $d$ minimo; indicare la complessità in tempo e memoria.

\textbf{Dove trovare la risposta}: richiede pseudocodice con tre cicli annidati; vedi il capitolo \emph{Programmazione Orientata agli Oggetti (e Fondamenti)}, sezione sulla complessità. Complessità: tempo $O(W \cdot H \cdot D)$, memoria $O(W \cdot H)$ per l'immagine di uscita.
